% Part: history
% Chapter: biographies 
% Section: georg-cantor

\documentclass[../../../include/open-logic-section]{subfiles}

\begin{document}

\olfileid[hi]{his}{bio}{can}

\olsection{गेऑर्ग कांटोर}

\olphoto{cantor-georg}{गेऑर्ग कांटोर}

गेऑर्ग कांटोर (\textsc{गे}-ऑर्ग \textsc{कान}-टोर) की एक आरंभिक जीवनी का
दावा था कि रूस के सेंट पीटर्सबर्ग जा रहे जहाज़ पर उनका जन्म हुआ और वे वहीं
मिले थे तथा उनके माता-पिता अज्ञात थे। किंतु यह सच नहीं है; हालाँकि उनका जन्म
1845 में सेंट पीटर्सबर्ग में ही हुआ था।

कांटोर ने 1867 में बर्लिन विश्वविद्यालय से गणित में डॉक्टरेट प्राप्त की। वे
समुच्चय सिद्धांत में अपने कार्य के लिए जाने जाते हैं और एक विशिष्ट अनुसंधान
विषय के रूप में समुच्चय सिद्धांत की स्थापना का श्रेय उन्हें दिया जाता है।
उन्होंने पहली बार सिद्ध किया कि भिन्न आकारों के अनंत समुच्चय होते हैं। उनके
सिद्धांतों, और विशेषतः उनके अनंतताओं के सिद्धांत, ने तत्कालीन गणितज्ञों में
बहुत वाद-विवाद उत्पन्न किया और उनका कार्य विवादास्पद रहा।

कांटोर की धार्मिक आस्थाएँ और उनका गणितीय कार्य अभिन्न रूप से जुड़े थे; उनका
दावा यहाँ तक था कि परिमितेतर संख्याओं का सिद्धांत स्वयं ईश्वर ने सीधे उन्हें
बताया था। जीवन के उत्तरार्ध में कांटोर मानसिक रोग से पीड़ित रहे। 1894 से और
अंतिम वर्षों की ओर अधिक बार उन्हें अस्पताल में भर्ती किया गया। उनके कार्य की
कड़ी आलोचना, जिसमें गणितज्ञ लियोपोल्ड क्रोनेकर से मतभेद भी था, अवसाद और गणित
में रुचि की कमी का कारण बनी। अवसाद के दौरों में कांटोर दर्शन और साहित्य की
ओर मुड़ते और उन्होंने यह सिद्धांत भी प्रकाशित किया कि शेक्सपियर के नाटकों के
लेखक फ़्रांसिस बेकन थे।

कांटोर का निधन 6 जनवरी 1918 को हाले के एक आरोग्याश्रम में हुआ।

\begin{reading} 
कांटोर की पूर्ण जीवनियों के लिए \citet{Dauben1990} और
\citet{Grattan-Guinness1971} देखें। कांटोर के आमूल विचारों का वर्णन बीबीसी
रेडियो 4 के कार्यक्रम \emph{A Brief History of Mathematics}
\citep{Sautoy2014} में भी है। यदि आप कांटोर के सिद्धांतों को रैप के रूप में
सुनना चाहें तो \citet{Rose2012} देखें।
\end{reading}

\end{document}
